\section{Conclusion}
\label{sec:conclusion}

This paper has presented a comprehensive empirical comparison of Classical Particle Swarm Optimization and Quantum-behaved Particle Swarm Optimization, implemented within a modular, open-source Python framework that supports three distinct quantum randomness generation modes. Through systematic evaluation on six standard benchmark functions across dimensionalities ranging from 2 to 20, with 30 independent runs per configuration and rigorous non-parametric statistical validation, we have established a nuanced picture of the conditions under which each algorithm excels.

Our findings demonstrate that QPSO with the weighted mean best position strategy achieves statistically significant improvements over Classical PSO on multimodal benchmark functions, particularly Rastrigin, Ackley, and the deceptive Schwefel function, where the quantum-inspired position update mechanism enables more effective exploration of rugged fitness landscapes. On the valley-shaped Rosenbrock function, the weighted mean best attractor provides a directional bias that facilitates navigation along the narrow ridge toward the global optimum, an advantage that becomes more pronounced as dimensionality increases. On smooth, unimodal landscapes such as the Sphere function, both algorithms perform comparably, confirming that the quantum-inspired exploration mechanism offers limited incremental benefit when the fitness landscape does not contain local optima to escape.

The comparison of the three quantum randomness modes---mathematical simulation, full Qiskit quantum circuit execution, and hybrid---reveals that solution quality is largely insensitive to the source of stochasticity, with the Math mode achieving performance within close proximity of the circuit-based modes across all benchmarks. The computational overhead of the Full mode, driven by circuit construction and classical simulation, is substantial, positioning the Math mode as the practical choice for current optimization applications while the circuit-based modes serve as a foundation for future deployment on real quantum processing units.

The principal contributions of this work are threefold. First, we provide a fair, statistically rigorous comparison that isolates algorithmic differences by employing paired random seeds, identical swarm sizes, and matched parameter schedules, thereby eliminating confounding factors that have affected prior comparative studies. Second, the modular framework architecture separates optimization logic from quantum randomness generation, benchmark function evaluation, and statistical analysis, enabling researchers to extend the system with new algorithms, benchmarks, or quantum backends with minimal effort. Third, the reproducible experimental pipeline---from raw optimization runs through convergence analysis, visualization, and hypothesis testing---establishes a methodological template for future comparative studies in quantum-inspired metaheuristics.

Several limitations of the present study should be acknowledged. All quantum circuit executions were performed using Qiskit's classical simulator rather than actual quantum hardware, and the randomness properties of simulated circuits differ from those of physical quantum measurements in ways that may affect optimizer behavior at scale. The benchmark suite, while spanning the major landscape categories, is limited to six functions and does not include the large-scale, rotated, and composite test suites used in competitions such as the IEEE Congress on Evolutionary Computation (CEC) benchmark series \cite{Liang2013}. Additionally, the parameter schedules for both algorithms were fixed based on literature recommendations rather than being tuned for each specific benchmark, which may not reflect the best achievable performance of either algorithm under optimized configurations.

These limitations suggest several promising directions for future research. Execution on real quantum hardware, now increasingly accessible through cloud-based quantum computing platforms, would allow investigation of whether genuine quantum randomness provides measurable advantages over classical pseudorandom and mathematically simulated alternatives. Expanding the benchmark suite to include the CEC 2014 or CEC 2022 competition functions, with their rotated, shifted, and composite landscapes, would provide a more comprehensive assessment of algorithmic robustness. Adaptive parameter tuning mechanisms, such as self-adaptive contraction--expansion schedules for QPSO or reinforcement learning-based parameter controllers, could further improve performance by tailoring the algorithm's exploration--exploitation balance to the characteristics of the landscape encountered during optimization. Finally, extending the framework to multi-objective optimization problems would broaden its applicability to real-world engineering design tasks where multiple competing objectives must be simultaneously optimized.

The complete framework, including all source code, benchmark implementations, experimental configurations, analysis scripts, and visualization tools, is available as an open-source Python package. We encourage the research community to use and extend this framework as a platform for reproducible experimentation in quantum-inspired optimization.
